\documentclass[../../main.tex]{subfiles}
\begin{document}

\begin{center}
  {\fontsize{48pt}{58pt}\selectfont \textbf{2000}\normalsize\textbf{年度}}
\end{center}

\begin{table}[htb]
  \centering
  \caption{出題テーマと難易度}
  \begin{tabular}{|c|p{9cm}|c|}
    \hline
    \textbf{問題番号} & \textbf{テーマ} & \textbf{難易度} \\
    \hline
    第1問 & 半円による光の反射と動く範囲 & \\
    \hline
    第2問 & 複素数の極形式と不等式の証明 & \\
    \hline
    第3問 & 三角柱の切断面の面積の範囲 & \\
    \hline
    第4問 & 指数関数のグラフの交点と面積の極限 & \\
    \hline
  \end{tabular}
\end{table}

\begin{table}[htb]
  \centering
  \caption{解答}
  \renewcommand{\arraystretch}{1.4}
  \begin{tabular}{|c|c|p{9cm}|}
    \hline
    \textbf{問題番号} & \textbf{小問} & \textbf{答え} \\
    \hline
    第1問 & (1) & $\dfrac{\pi}{4}<\theta<\dfrac{\pi}{3}$ \\
    \cline{2-3}
    & (2) & $P=\left(\dfrac{\sin\theta}{\sin5\theta},\ 0\right)$ \\
    \cline{2-3}
    & (3) & $-1<x\le-\dfrac{4}{5}$，$y=0$ \\
    \hline
    第2問 & (1) & $r>0$かつ$r<-\cos\theta$ \\
    \cline{2-3}
    & (2) & $|Z_n|^2=\dfrac{r^{2n+2}-2r^{n+1}\cos(n+1)\theta+1}{r^2-2r\cos\theta+1}$ \\
    \cline{2-3}
    & (3) & 証明問題 \\
    \hline
    第3問 & - & $\dfrac{3}{4}\le S\le\dfrac{\sqrt{17}}{4}$ \\
    \hline
    第4問 & (1) & 証明問題 \\
    \cline{2-3}
    & (2) & $\displaystyle\lim_{n\to\infty}a_n=0$，$\displaystyle\lim_{n\to\infty}na_n=\log2$ \\
    \cline{2-3}
    & (3) & $2\log2-1$ \\
    \hline
  \end{tabular}
\end{table}

\paragraph{試験形態}
4問．試験時間は要確認．

\end{document}
